% Options for packages loaded elsewhere
\PassOptionsToPackage{unicode}{hyperref}
\PassOptionsToPackage{hyphens}{url}
\documentclass[
]{article}
\usepackage{xcolor}
\usepackage[margin=1in]{geometry}
\usepackage{amsmath,amssymb}
\setcounter{secnumdepth}{-\maxdimen} % remove section numbering
\usepackage{iftex}
\ifPDFTeX
  \usepackage[T1]{fontenc}
  \usepackage[utf8]{inputenc}
  \usepackage{textcomp} % provide euro and other symbols
\else % if luatex or xetex
  \usepackage{unicode-math} % this also loads fontspec
  \defaultfontfeatures{Scale=MatchLowercase}
  \defaultfontfeatures[\rmfamily]{Ligatures=TeX,Scale=1}
\fi
\usepackage{lmodern}
\ifPDFTeX\else
  % xetex/luatex font selection
\fi
% Use upquote if available, for straight quotes in verbatim environments
\IfFileExists{upquote.sty}{\usepackage{upquote}}{}
\IfFileExists{microtype.sty}{% use microtype if available
  \usepackage[]{microtype}
  \UseMicrotypeSet[protrusion]{basicmath} % disable protrusion for tt fonts
}{}
\makeatletter
\@ifundefined{KOMAClassName}{% if non-KOMA class
  \IfFileExists{parskip.sty}{%
    \usepackage{parskip}
  }{% else
    \setlength{\parindent}{0pt}
    \setlength{\parskip}{6pt plus 2pt minus 1pt}}
}{% if KOMA class
  \KOMAoptions{parskip=half}}
\makeatother
\usepackage{color}
\usepackage{fancyvrb}
\newcommand{\VerbBar}{|}
\newcommand{\VERB}{\Verb[commandchars=\\\{\}]}
\DefineVerbatimEnvironment{Highlighting}{Verbatim}{commandchars=\\\{\}}
% Add ',fontsize=\small' for more characters per line
\usepackage{framed}
\definecolor{shadecolor}{RGB}{248,248,248}
\newenvironment{Shaded}{\begin{snugshade}}{\end{snugshade}}
\newcommand{\AlertTok}[1]{\textcolor[rgb]{0.94,0.16,0.16}{#1}}
\newcommand{\AnnotationTok}[1]{\textcolor[rgb]{0.56,0.35,0.01}{\textbf{\textit{#1}}}}
\newcommand{\AttributeTok}[1]{\textcolor[rgb]{0.13,0.29,0.53}{#1}}
\newcommand{\BaseNTok}[1]{\textcolor[rgb]{0.00,0.00,0.81}{#1}}
\newcommand{\BuiltInTok}[1]{#1}
\newcommand{\CharTok}[1]{\textcolor[rgb]{0.31,0.60,0.02}{#1}}
\newcommand{\CommentTok}[1]{\textcolor[rgb]{0.56,0.35,0.01}{\textit{#1}}}
\newcommand{\CommentVarTok}[1]{\textcolor[rgb]{0.56,0.35,0.01}{\textbf{\textit{#1}}}}
\newcommand{\ConstantTok}[1]{\textcolor[rgb]{0.56,0.35,0.01}{#1}}
\newcommand{\ControlFlowTok}[1]{\textcolor[rgb]{0.13,0.29,0.53}{\textbf{#1}}}
\newcommand{\DataTypeTok}[1]{\textcolor[rgb]{0.13,0.29,0.53}{#1}}
\newcommand{\DecValTok}[1]{\textcolor[rgb]{0.00,0.00,0.81}{#1}}
\newcommand{\DocumentationTok}[1]{\textcolor[rgb]{0.56,0.35,0.01}{\textbf{\textit{#1}}}}
\newcommand{\ErrorTok}[1]{\textcolor[rgb]{0.64,0.00,0.00}{\textbf{#1}}}
\newcommand{\ExtensionTok}[1]{#1}
\newcommand{\FloatTok}[1]{\textcolor[rgb]{0.00,0.00,0.81}{#1}}
\newcommand{\FunctionTok}[1]{\textcolor[rgb]{0.13,0.29,0.53}{\textbf{#1}}}
\newcommand{\ImportTok}[1]{#1}
\newcommand{\InformationTok}[1]{\textcolor[rgb]{0.56,0.35,0.01}{\textbf{\textit{#1}}}}
\newcommand{\KeywordTok}[1]{\textcolor[rgb]{0.13,0.29,0.53}{\textbf{#1}}}
\newcommand{\NormalTok}[1]{#1}
\newcommand{\OperatorTok}[1]{\textcolor[rgb]{0.81,0.36,0.00}{\textbf{#1}}}
\newcommand{\OtherTok}[1]{\textcolor[rgb]{0.56,0.35,0.01}{#1}}
\newcommand{\PreprocessorTok}[1]{\textcolor[rgb]{0.56,0.35,0.01}{\textit{#1}}}
\newcommand{\RegionMarkerTok}[1]{#1}
\newcommand{\SpecialCharTok}[1]{\textcolor[rgb]{0.81,0.36,0.00}{\textbf{#1}}}
\newcommand{\SpecialStringTok}[1]{\textcolor[rgb]{0.31,0.60,0.02}{#1}}
\newcommand{\StringTok}[1]{\textcolor[rgb]{0.31,0.60,0.02}{#1}}
\newcommand{\VariableTok}[1]{\textcolor[rgb]{0.00,0.00,0.00}{#1}}
\newcommand{\VerbatimStringTok}[1]{\textcolor[rgb]{0.31,0.60,0.02}{#1}}
\newcommand{\WarningTok}[1]{\textcolor[rgb]{0.56,0.35,0.01}{\textbf{\textit{#1}}}}
\usepackage{graphicx}
\makeatletter
\newsavebox\pandoc@box
\newcommand*\pandocbounded[1]{% scales image to fit in text height/width
  \sbox\pandoc@box{#1}%
  \Gscale@div\@tempa{\textheight}{\dimexpr\ht\pandoc@box+\dp\pandoc@box\relax}%
  \Gscale@div\@tempb{\linewidth}{\wd\pandoc@box}%
  \ifdim\@tempb\p@<\@tempa\p@\let\@tempa\@tempb\fi% select the smaller of both
  \ifdim\@tempa\p@<\p@\scalebox{\@tempa}{\usebox\pandoc@box}%
  \else\usebox{\pandoc@box}%
  \fi%
}
% Set default figure placement to htbp
\def\fps@figure{htbp}
\makeatother
\setlength{\emergencystretch}{3em} % prevent overfull lines
\providecommand{\tightlist}{%
  \setlength{\itemsep}{0pt}\setlength{\parskip}{0pt}}
\usepackage{bookmark}
\IfFileExists{xurl.sty}{\usepackage{xurl}}{} % add URL line breaks if available
\urlstyle{same}
\hypersetup{
  pdftitle={Sessão 3: Exercícios Práticos},
  pdfauthor={Os vossos formadores},
  hidelinks,
  pdfcreator={LaTeX via pandoc}}

\title{Sessão 3: Exercícios Práticos}
\author{Os vossos formadores}
\date{6 de Julho de 2025}

\begin{document}
\maketitle

{
\setcounter{tocdepth}{3}
\tableofcontents
}
\section{Exercícios Práticos: dplyr, NAs e Novas
Variáveis}\label{exercuxedcios-pruxe1ticos-dplyr-nas-e-novas-variuxe1veis}

Olá! Nesta sessão de exercícios, vamos consolidar os vossos
conhecimentos sobre manipulação de dados com dplyr, tratamento de
valores em falta e criação de novas variáveis. Vamos trabalhar com um
dataset simulado de ``Dados de Clientes''.

\textbf{Preparação:} Criar o Dataset de Exemplo Primeiro, vamos criar um
ficheiro CSV de exemplo com alguns valores em falta para os exercícios.
Certifiquem-se de que o vosso diretório de trabalho está correto ou que
o RStudio Project está ativo.

\begin{Shaded}
\begin{Highlighting}[]
\CommentTok{\# Criar um data frame de exemplo com NAs}
\NormalTok{dados\_clientes }\OtherTok{\textless{}{-}} \FunctionTok{data.frame}\NormalTok{(}
  \AttributeTok{ID\_Cliente =} \DecValTok{1001}\SpecialCharTok{:}\DecValTok{1010}\NormalTok{,}
  \AttributeTok{Nome =} \FunctionTok{c}\NormalTok{(}\StringTok{"Alice"}\NormalTok{, }\StringTok{"Bernardo"}\NormalTok{, }\StringTok{"Carla"}\NormalTok{, }\StringTok{"Diogo"}\NormalTok{, }\StringTok{"Eva"}\NormalTok{, }\StringTok{"Filipe"}\NormalTok{, }\StringTok{"Gabriela"}\NormalTok{, }\StringTok{"Hugo"}\NormalTok{, }\StringTok{"Inês"}\NormalTok{, }\StringTok{"Joana"}\NormalTok{),}
  \AttributeTok{Idade =} \FunctionTok{c}\NormalTok{(}\DecValTok{25}\NormalTok{, }\DecValTok{30}\NormalTok{, }\ConstantTok{NA}\NormalTok{, }\DecValTok{22}\NormalTok{, }\DecValTok{35}\NormalTok{, }\DecValTok{28}\NormalTok{, }\ConstantTok{NA}\NormalTok{, }\DecValTok{40}\NormalTok{, }\DecValTok{27}\NormalTok{, }\DecValTok{33}\NormalTok{),}
  \AttributeTok{Rendimento\_Mensal =} \FunctionTok{c}\NormalTok{(}\DecValTok{1500}\NormalTok{, }\DecValTok{2200}\NormalTok{, }\DecValTok{1800}\NormalTok{, }\DecValTok{1400}\NormalTok{, }\ConstantTok{NA}\NormalTok{, }\DecValTok{2500}\NormalTok{, }\DecValTok{1900}\NormalTok{, }\DecValTok{3000}\NormalTok{, }\ConstantTok{NA}\NormalTok{, }\DecValTok{2100}\NormalTok{),}
  \AttributeTok{Genero =} \FunctionTok{c}\NormalTok{(}\StringTok{"F"}\NormalTok{, }\StringTok{"M"}\NormalTok{, }\StringTok{"F"}\NormalTok{, }\StringTok{"M"}\NormalTok{, }\StringTok{"F"}\NormalTok{, }\StringTok{"M"}\NormalTok{, }\StringTok{"F"}\NormalTok{, }\StringTok{"M"}\NormalTok{, }\StringTok{"F"}\NormalTok{, }\StringTok{"M"}\NormalTok{),}
  \AttributeTok{Cidade =} \FunctionTok{c}\NormalTok{(}\StringTok{"Lisboa"}\NormalTok{, }\StringTok{"Porto"}\NormalTok{, }\StringTok{"Coimbra"}\NormalTok{, }\StringTok{"Braga"}\NormalTok{, }\StringTok{"Lisboa"}\NormalTok{, }\StringTok{"Porto"}\NormalTok{, }\StringTok{"Coimbra"}\NormalTok{, }\StringTok{"Braga"}\NormalTok{, }\StringTok{"Lisboa"}\NormalTok{, }\StringTok{"Porto"}\NormalTok{),}
  \AttributeTok{Compras\_Ultimos\_6\_Meses =} \FunctionTok{c}\NormalTok{(}\DecValTok{5}\NormalTok{, }\DecValTok{8}\NormalTok{, }\DecValTok{3}\NormalTok{, }\DecValTok{10}\NormalTok{, }\DecValTok{6}\NormalTok{, }\DecValTok{7}\NormalTok{, }\DecValTok{4}\NormalTok{, }\DecValTok{12}\NormalTok{, }\DecValTok{9}\NormalTok{, }\DecValTok{5}\NormalTok{)}
\NormalTok{)}

\CommentTok{\# Introduzir alguns NAs adicionais para testar}
\NormalTok{dados\_clientes}\SpecialCharTok{$}\NormalTok{Rendimento\_Mensal[}\FunctionTok{c}\NormalTok{(}\DecValTok{5}\NormalTok{, }\DecValTok{9}\NormalTok{)] }\OtherTok{\textless{}{-}} \ConstantTok{NA}
\NormalTok{dados\_clientes}\SpecialCharTok{$}\NormalTok{Idade[}\DecValTok{3}\NormalTok{] }\OtherTok{\textless{}{-}} \ConstantTok{NA}

\CommentTok{\# Guardar o data frame num ficheiro CSV}
\FunctionTok{write.csv}\NormalTok{(dados\_clientes, }\StringTok{"dados\_clientes.csv"}\NormalTok{, }\AttributeTok{row.names =} \ConstantTok{FALSE}\NormalTok{)}

\FunctionTok{print}\NormalTok{(}\StringTok{"Ficheiro \textquotesingle{}dados\_clientes.csv\textquotesingle{} criado na vossa pasta de trabalho."}\NormalTok{)}
\end{Highlighting}
\end{Shaded}

\begin{verbatim}
## [1] "Ficheiro 'dados_clientes.csv' criado na vossa pasta de trabalho."
\end{verbatim}

\begin{Shaded}
\begin{Highlighting}[]
\CommentTok{\# Carregar o pacote dplyr}
\FunctionTok{library}\NormalTok{(dplyr)}
\end{Highlighting}
\end{Shaded}

\begin{verbatim}
## 
## Attaching package: 'dplyr'
\end{verbatim}

\begin{verbatim}
## The following objects are masked from 'package:stats':
## 
##     filter, lag
\end{verbatim}

\begin{verbatim}
## The following objects are masked from 'package:base':
## 
##     intersect, setdiff, setequal, union
\end{verbatim}

\begin{Shaded}
\begin{Highlighting}[]
\CommentTok{\# Carregar o pacote tidyr (para drop\_na e replace\_na)}
\FunctionTok{library}\NormalTok{(tidyr)}

\CommentTok{\# Importar o dataset}
\NormalTok{meus\_clientes }\OtherTok{\textless{}{-}} \FunctionTok{read.csv}\NormalTok{(}\StringTok{"dados\_clientes.csv"}\NormalTok{)}
\FunctionTok{head}\NormalTok{(meus\_clientes)}
\end{Highlighting}
\end{Shaded}

\begin{verbatim}
##   ID_Cliente     Nome Idade Rendimento_Mensal Genero  Cidade
## 1       1001    Alice    25              1500      F  Lisboa
## 2       1002 Bernardo    30              2200      M   Porto
## 3       1003    Carla    NA              1800      F Coimbra
## 4       1004    Diogo    22              1400      M   Braga
## 5       1005      Eva    35                NA      F  Lisboa
## 6       1006   Filipe    28              2500      M   Porto
##   Compras_Ultimos_6_Meses
## 1                       5
## 2                       8
## 3                       3
## 4                      10
## 5                       6
## 6                       7
\end{verbatim}

\subsection{\texorpdfstring{\textbf{Exercício 1:} Exploração e
Identificação de
NAs}{Exercício 1: Exploração e Identificação de NAs}}\label{exercuxedcio-1-explorauxe7uxe3o-e-identificauxe7uxe3o-de-nas}

\begin{enumerate}
\def\labelenumi{\arabic{enumi}.}
\tightlist
\item
  \textbf{Explorar:} Obtenham um summary() e str() do meus\_clientes.
\item
  \textbf{Contar NAs:} Quantos valores em falta existem em cada coluna
  (Idade e Rendimento\_Mensal)?
\item
  \textbf{Linhas Completas:} Quantas linhas no data frame não têm NAs?
\end{enumerate}

\begin{Shaded}
\begin{Highlighting}[]
\CommentTok{\# Escrevam o vosso código aqui para o Exercício 1}
\end{Highlighting}
\end{Shaded}

\subsection{\texorpdfstring{\textbf{Exercício 2:} Tratamento de Valores
em
Falta}{Exercício 2: Tratamento de Valores em Falta}}\label{exercuxedcio-2-tratamento-de-valores-em-falta}

\begin{enumerate}
\def\labelenumi{\arabic{enumi}.}
\tightlist
\item
  \textbf{Remover NAs:} Criem um novo data frame chamado
  clientes\_sem\_na que remova todas as linhas com valores em falta.
  Quantas linhas restaram?
\item
  \textbf{Imputar Idade:} Criem um novo data frame chamado
  clientes\_imputado\_idade onde os NAs na coluna Idade são substituídos
  pela média da Idade (ignorando os NAs para o cálculo da média).
\item
  \textbf{Imputar Rendimento:} No clientes\_imputado\_idade, substituam
  os NAs na coluna Rendimento\_Mensal pela mediana do Rendimento\_Mensal
  (ignorando os NAs para o cálculo da mediana). Chame o resultado final
  de clientes\_limpo.
\end{enumerate}

\begin{Shaded}
\begin{Highlighting}[]
\CommentTok{\# Escrevam o vosso código aqui para o Exercício 2}
\end{Highlighting}
\end{Shaded}

\subsection{\texorpdfstring{\textbf{Exercício 3:} Criação de Novas
Variáveis}{Exercício 3: Criação de Novas Variáveis}}\label{exercuxedcio-3-criauxe7uxe3o-de-novas-variuxe1veis}

Continuem a trabalhar com o clientes\_limpo (o data frame sem NAs e com
imputação).

\begin{enumerate}
\def\labelenumi{\arabic{enumi}.}
\tightlist
\item
  \textbf{Rendimento Anual:} Criem uma nova coluna chamada
  Rendimento\_Anual que seja o Rendimento\_Mensal multiplicado por 12.
\item
  \textbf{Categoria de Idade:} Criem uma nova coluna chamada
  Categoria\_Idade com base na Idade:

  \begin{itemize}
  \tightlist
  \item
    ``Jovem'' se Idade \textless{} 25
  \item
    ``Adulto'' se Idade \textgreater= 25 e Idade \textless{} 35
  \item
    ``Senior'' se Idade \textgreater= 35
  \end{itemize}
\item
  \textbf{Frequência de Compra:} Criem uma nova coluna chamada
  Frequencia\_Compra com base em Compras\_Ultimos\_6\_Meses:

  \begin{itemize}
  \tightlist
  \item
    ``Baixa'' se Compras\_Ultimos\_6\_Meses \textless= 5
  \item
    ``Média'' se Compras\_Ultimos\_6\_Meses \textgreater{} 5 e
    Compras\_Ultimos\_6\_Meses \textless= 8
  \item
    ``Alta'' se Compras\_Ultimos\_6\_Meses \textgreater{} 8
  \end{itemize}
\end{enumerate}

\begin{Shaded}
\begin{Highlighting}[]
\CommentTok{\# Escrevam o vosso código aqui para o Exercício 3}
\end{Highlighting}
\end{Shaded}

\subsection{\texorpdfstring{\textbf{Exercício 4:} Agrupamento e Resumo
com
dplyr}{Exercício 4: Agrupamento e Resumo com dplyr}}\label{exercuxedcio-4-agrupamento-e-resumo-com-dplyr}

Continuem a trabalhar com o clientes\_limpo (com as novas variáveis
criadas).

\begin{enumerate}
\def\labelenumi{\arabic{enumi}.}
\tightlist
\item
  Média de Idade por Género: Calculem a média da Idade para cada Genero.
\item
  Rendimento Médio por Cidade: Calculem o Rendimento\_Mensal médio e o
  número de clientes (n()) para cada Cidade.
\item
  Resumo por Categoria de Idade e Género: Calculem o Rendimento\_Anual
  médio e o número de clientes (n()) para cada combinação de
  Categoria\_Idade e Genero. Ordenem os resultados pelo
  Rendimento\_Anual médio em ordem decrescente.
\end{enumerate}

\begin{Shaded}
\begin{Highlighting}[]
\CommentTok{\# Escrevam o vosso código aqui para o Exercício 4}
\end{Highlighting}
\end{Shaded}

\textbf{Parabéns!} Completaram os exercícios da Sessão 3!

\end{document}
